\documentclass[11pt, letterpaper]{article}
\input{../../homework.tex}
\usepackage{titlepageBU}
\title{Turn in 6}
\courseID{MA 542}
\courseSection{A1}
\begin{document}
\makeproblem
\section*{Problem 1}
Since $\sqrt[5]{7} ,\zeta \sqrt[5]{7} \in E$, we have
\[
	\frac{\zeta \sqrt[5]{7} }{\sqrt[5]{7} } =\zeta \in E
.\]
Therefore, $\mathbb{Q} \subseteq \mathbb{Q} (\zeta )\subseteq E$ is a chain of field extensios, and thus
\[
	\left[ E:\mathbb{Q} (\zeta ) \right] \left[ \mathbb{Q} (\zeta ):\mathbb{Q}  \right] .
\]
It suffices to compute these two smaller computations. Note that $\mathbb{Q} (\zeta )$ is generated by the set $Z\coloneqq \left\{ \zeta ^i\mid 0\leq i\leq 4,i\in\mathbb{N}  \right\} $. Note also that $\zeta $ manifestly satisfies the equation $\zeta ^5-1=0$, and that
\begin{align*}
	\left( \zeta -1 \right) \left( \zeta ^4+\zeta ^3+\zeta ^2+\zeta ^1+1 \right) &=\zeta ^5+\zeta ^4+\zeta ^3+\zeta ^2+\zeta -\zeta ^4-\zeta ^3-\zeta ^2-\zeta -1\\
										     &=\zeta ^5-1\\
										     &=0
.\end{align*}
Since $\zeta $ is a primitive fifth root of unity, $\zeta \neq 1$, and thus $\zeta -1\neq 0$. Dividing by this, we find that
\[
	\zeta ^4+\zeta ^3+\zeta ^2+\zeta +1=0
.\]
Therefore, this set is not linearly independent. Moreover, since $\zeta ^5=1$, we have
\[
	\left( \zeta ^2 \right) ^4+\left( \zeta ^2 \right) ^3+\left( \zeta ^2 \right) ^2+\left( \zeta ^2 \right) +1=\zeta ^8+\zeta ^6+\zeta ^4+\zeta ^2+1=\zeta ^3+\zeta +\zeta ^4+\zeta ^2+1=0,
\]
\[
	\left( \zeta ^3 \right) ^4+\left( \zeta ^3 \right) ^3+\left( \zeta ^3 \right) ^2+\left( \zeta ^3 \right) +1=\zeta ^{12} +\zeta ^9+\zeta ^6+\zeta ^3+1=\zeta ^2+\zeta ^4+\zeta +\zeta ^3+1=0,
\]
\[
	\left( \zeta ^4 \right) ^4+\left( \zeta ^4 \right) ^3+\left( \zeta ^4 \right) ^2+\left( \zeta ^4 \right) +1=\zeta ^{16} +\zeta ^{12} +\zeta ^8+\zeta ^4+1=\zeta +\zeta ^2+\zeta ^3+\zeta ^4+1=0
.\]
Therefore, $\left\{ \zeta ,\zeta ^2,\zeta ^3,\zeta ^4 \right\} $ form a set of four distinct roots of this polynomial. Since a polynomial in $\mathbb{Q} $ has at most 4 roots, these must be all roots of the polynomial. Since none of these roots are in $\mathbb{Q} $, we find that this polynomial is irreducible over $\mathbb{Q} $, and thus is the minimal polynomial of $\zeta $. Since the polynomial has degree 4, $\left[ \mathbb{Q}(\zeta ):\mathbb{Q}  \right] =4$.

Now consider $E$. Note that $x^5-7$ is clearly the minimal polynomial of $\sqrt[5]{7} $ over $\mathbb{Q} $. However, note that $x^5-7$ is also clearly irreducible over $\mathbb{Q} (\zeta )$ as well, and thus is the minimal polynomial over $\mathbb{Q} (\zeta )$. Since $E$ is the minimal field containing $\zeta $ and $\sqrt[5]{7} $, we find that $\left[ E:\mathbb{Q} (\zeta ) \right] =\operatorname{deg} (x^5-7)=5$. Therefore,
\[
	\left[ E:\mathbb{Q}  \right] =5\cdot 4=20
.\]
\section*{Problem 2}
Let $K_6=\mathbb{Q} (\zeta )$. The proof follows from our proof of the above problem.

Let $K_i=\mathbb{Q} (\zeta^{i}\sqrt[5]{7} )$ for $1\leq i\leq 5$. Since each $\zeta ^i \sqrt[5]{7} $ is a root of $x^5-7$, which is by Eisenstein's critera irreducible over $\mathbb{Q} $, we find that it is the minimal polynomial for each root over $\mathbb{Q} $, and thus $\left[ \mathbb{Q} (\zeta ^i \sqrt[5]{7} ):\mathbb{Q}  \right] =\operatorname{deg} \left( x^5-7 \right) =5$.
\section*{Problem 3}
Since each $K_i$ for $1\leq i\leq 5$ does not contain all roots of $x^5-7$, the polynomial does not split in $K_i$. Since this is the minimal polynomial for each of these roots, it follows that no polynomial that doesn't split in $\mathbb{Q} $ can split in these fields. Therefore, these fields are not splitting fields.

The field $\mathbb{Q} (\zeta )$ is a splitting field for $x^4+x^3+x^2+x+1$, as shown previously. Therefore, it is a splitting field.
\section*{Problem 4}
Since the degree of $E$ is 20, and since $E$ is clearly Galois as it has characteristic 0 and $x^5-7$ splits in it, we have $\left[ E:\mathbb{Q}  \right] =\left| \operatorname{Gal} (E/\mathbb{Q} ) \right| =20$. Since $\sigma ,\tau $ are clearly well-defined as automorphisms by definition, they are elements of $\operatorname{Gal} (E/\mathbb{Q} )$. It suffices to show $|\left\langle \sigma ,\tau  \right\rangle| =20$.

Note that $\left| \sigma  \right| =5$ trivially, as $\sigma $ fixes all elements other than powers of $\sqrt[5]{7} $, and
\[
	\sigma ^5(\sqrt[5]{7} )=\zeta ^5\sqrt[5]{7} =\sqrt[5]{7} 
.\]
Since this generates all other powers of $\sqrt[5]{7} $, it is easily seen that $\left| \sigma  \right| =5$. Similarly, $\tau $ fixes $\sqrt[5]{7} $, but it's action on $\zeta $ gives it order 4:
\[
	\tau (\zeta )=\zeta ^2,\qquad \tau ^2(\zeta )=\left( \zeta ^2 \right) ^2=\zeta ^2,
\]
\[
	\tau ^3(\zeta )=\left( \zeta ^4 \right) ^2=\zeta ^8,\qquad \tau ^4(\zeta )=\left( \zeta ^8 \right) ^2=\zeta ^{16} =\zeta 
.\]
Next, note that
\[
	\left( \tau \sigma \tau ^{-1}  \right) (\zeta )=\tau \tau ^{-1} (\zeta )=\zeta ,
\]
since $\sigma $ fixes $\zeta $, and also note that
\[
	(\tau \sigma \tau ^{-1} )\left( \sqrt[5]{7}  \right) =(\tau \sigma )\left( \sqrt[5]{7}  \right) =\zeta ^2\sqrt[5]{7} ,
\]
since $\tau $ fixes the radical, and thus so must it's inverse. Also notice that
\[
	\sigma ^2(\zeta )=\zeta ,\qquad \sigma ^2\left( \sqrt[5]{7}  \right) =\zeta ^2\sqrt[5]{7} 
.\]
Therefore, $\sigma ^2=\tau \sigma \tau ^{-1} $. Note that $\sigma ^2$ generates $\left\langle \sigma  \right\rangle $, since $2$ doesn't divide 5. Therefore,
\[
	\tau \left\langle \sigma  \right\rangle \tau ^{-1} =\left\langle \sigma  \right\rangle ,
\]
and thus we immediately have $\tau ^n \left\langle \sigma  \right\rangle \tau ^{-n} =\left\langle \sigma  \right\rangle $. It follows that $\left\langle \sigma  \right\rangle $ is normal in $\left\langle \sigma ,\tau  \right\rangle $, and thus
\[
	\left\langle \tau  \right\rangle \left\langle \sigma  \right\rangle =\bigcup_{\tau^i \in \left\langle \tau  \right\rangle } \tau ^i \left\langle \sigma  \right\rangle =\bigcup_{\tau ^i\in \left\langle \tau  \right\rangle } \left\langle \sigma  \right\rangle \tau ^{i} =\left\langle \sigma  \right\rangle \left\langle \tau  \right\rangle 
.\]
By a theorem we proved and used extensively last semester, $\left\langle \sigma  \right\rangle \left\langle \tau  \right\rangle $ is a subgroup of $\left\langle \sigma ,\tau  \right\rangle $. Also note that since $\tau $ fixes the element $\sigma $ acts on and vice versa, $\left\langle \tau  \right\rangle \cap \left\langle \sigma  \right\rangle =1$ trivially. Therefore, $\left\langle \tau  \right\rangle \left\langle \sigma  \right\rangle $ is an internal direct product, and thus
\[
	\left\langle \tau  \right\rangle \left\langle \sigma  \right\rangle \cong \left\langle \tau  \right\rangle \times \left\langle \sigma  \right\rangle 
.\]
This direct product of groups has order $\left| \left\langle \tau  \right\rangle  \right| \cdot \left| \left\langle \sigma  \right\rangle  \right| =5\cdot 4=20$, so we have that $\left\langle \tau ,\sigma  \right\rangle =\left\langle \tau  \right\rangle \left\langle \sigma  \right\rangle $ since they can both have no more than order 20, since they are subgroups of the Galois group which has order 20. We thus have $\left| \left\langle \tau ,\sigma  \right\rangle  \right| =20$, and thus $\left\langle \tau ,\sigma  \right\rangle =\operatorname{Gal} (E/\mathbb{Q} )$.

\end{document}
